\documentclass{article}
\PassOptionsToPackage{numbers,sort&compress}{natbib}
\usepackage[dblblindworkshop]{neurips_2026}
\workshoptitle{Foundation Models for Temporal Systems: From Forecasting to World Modeling}
\usepackage[utf8]{inputenc}
\usepackage[T1]{fontenc}
\usepackage{amsmath,amssymb,graphicx,booktabs,longtable,array,calc}
\usepackage{microtype,url}
\usepackage[hidelinks,unicode]{hyperref}
\providecommand{\tightlist}{\setlength{\itemsep}{0pt}\setlength{\parskip}{0pt}}
\providecommand{\pandocbounded}[1]{#1}
\setlength{\emergencystretch}{1em}
\title{Hidden-Feature Matching and Left-Right Movement: Different Outcomes in a Gait JEPA}
\author{Anonymous Author(s)}
\hypersetup{pdftitle={Hidden-Feature Matching and Left-Right Movement: Different Outcomes in a Gait JEPA},pdfauthor={},pdfcreator={LaTeX}}
\begin{document}
\maketitle
\begin{abstract}
We test whether learning to predict hidden skeleton features helps
estimate a left--right difference in movement speed during walking. A
small skeleton Joint-Embedding Predictive Architecture (JEPA) learns
without movement labels; a separate regression then estimates a signed
speed contrast from its frozen features. In source-video
cross-validation on 625 clips from 93 videos, trained predictors more
consistently favor hidden features from the correct clip over features
from another video. With the same expanded feature summary and
regression procedure, however, mean source-balanced \(R^2\) falls from
0.22 at initialization to 0.10--0.11 after training. Motion-weighted and
connected-region masks leave their movement-readout benefits uncertain.
The comparison identifies a gap between hidden-feature matching and
access to a specified observable of an evolving gait system, within this
cohort, training recipe and readout. It motivates evaluating observable
movement alongside latent prediction in representations intended for
temporal modeling.
\end{abstract}
\section{Introduction}\label{introduction}

Walking involves coupled body landmarks whose positions change over
time. Their trajectories carry information about movement that a single
pose cannot supply, while video-derived coordinates capture only part of
the physical process: camera geometry and pose errors affect them, and
forces and contact remain unobserved. Gait rhythm, phase and
coordination make this an appealing temporal system to study. Temporal
world modeling aims to represent an evolving system well enough to
predict its observable behavior; our experiment examines access to one
movement measurement within that broader objective.

We use a small S-JEPA \citep{sjepa} to predict features at hidden
positions in a skeleton sequence, then ask a separate regression to
estimate a left--right contrast of median coordinate speeds. The central
null is no improvement from JEPA training over matched initial weights
on held-out source videos. This null also motivates the registered
reflection test; the mask comparisons and expanded readout were
developed on the same cohort and remain exploratory. By comparing the
training task with a measurement of the observed movement, the study
offers focused evidence for evaluating representations used in temporal
foundation models.

\section{Data, prediction task and
evaluation}\label{data-prediction-task-and-evaluation}

\textbf{Measuring movement.} Quality checks retain 625 GAVD clips from
93 source videos \citep{gavd}. Archives contain \(T_i\times33\times4\)
arrays of three estimated coordinates and visibility, plus frame numbers
and frame rate. The coordinates are image normalized with inferred
depth, rather than calibrated metric positions. We center observed
coordinates at the pelvis and divide by body width, keeping original
timestamps and leaving gaps unfilled. Five pairs---shoulders, knees,
ankles, heels and foot tips---each need at least eight transitions with
both sides observed at both ends. For left and right median coordinate
speeds \(m_{L,k},m_{R,k}\), the signed contrast is

\[y=\frac{1}{5}\sum_{k=1}^{5}\frac{m_{L,k}-m_{R,k}}{m_{L,k}+m_{R,k}+10^{-8}}.\]

Reflection \(M\) flips the centered horizontal coordinate and exchanges
anatomical sides, including validity flags, so \(y(Mx)=-y(x)\). Hips
establish the pelvis but are excluded from the sum because centering
equalizes their speed magnitudes. The contrast measures relative median
speed. It omits phase, stride timing and loading, and remains unchanged
when a trajectory and its sampling intervals are reversed. Its recovery
alone therefore provides limited evidence about temporal dynamics.

\textbf{Preparing encoder inputs.} A landmark is observed when
visibility is at least 0.45 and its coordinates are finite. A separate
path fills interior gaps of up to four samples and treats those
positions as valid. After pelvis centering and body-width normalization,
coordinates and validity are linearly resized by sequence index to 64
positions; original time intervals are discarded. Resized validity must
reach 0.999, and invalid coordinates are zero. Each joint's
four-position patch projects 12 coordinates to 96 features, producing
\(16\times33\) tokens. All four positions must be valid for a patch to
be used. Anatomical identities and prepared positions remain, but
original duration and cadence are not retained; trajectory-derived
quantities can also change.

\begin{figure}[htbp]
\centering
\includegraphics[width=\linewidth]{assets/v8/figures/training_pipeline_compact.pdf}
\caption{Completed model. Poses and tokens are schematic: skeleton colors mark anatomical left (teal) and right (purple), with an open circle for a missing landmark. Tokens are visible (teal), hidden (orange) or missing (gray). Scattered/motion-weighted and connected-region masks may leave context on both sides. Full clips reach the teacher and regularizer; target y supervises the separate readout.}
\end{figure}

\textbf{Learning hidden features.} The online encoder zeroes hidden
patch content before adding joint and position embeddings. The predictor
inserts learned mask tokens at target positions and estimates the
teacher's features there. Cross-entropy compares distributions over
feature channels at hidden valid positions, averaging within clips and
then across clips. The teacher receives the full prepared clip; its
weights follow an exponential moving average of the online weights
without gradients. A VICReg regularizer \citep{vicreg}, weighted 0.05,
also uses full clips: two geometric views pass through the online
encoder, a twelve-joint gait pool including hips, and a projector.
Movement targets and gait-category labels never enter these losses.

The grid contains 125 encoders: three motion-mask conditions and two
region conditions across five folds and five seeds, each trained for
1,200 AdamW updates in CUDA BF16 and evaluated in FP32. Motion
conditions use random, MAMP-style \citep{mamp} or mixed
uniform/clipped-motion sampling over all 33 joints, with counts derived
from a twelve-joint gait pool. Region masks hide six connected landmarks
over eight of sixteen blocks. Each family has its own count-matched
random reference and shares initial weights, source draws, views and
update counts. Reflection augmentation and explicit reflection loss are
absent from this grid. Appendix A gives settings and their separate
evidence.

\textbf{Estimating movement from frozen features.} Five outer folds keep
every video's clips, overlapping windows, mirrors and augmented views
together wherever these derivatives are used. A fresh encoder fits only
outer-training sources, sampling videos equally. Its matched
initialization and final encoders are then frozen for ridge regression,
a linear readout with a coefficient penalty. Three inner source folds
select that penalty using outer-training labels only. Imputation and
scaling fit within each inner training split, then refit with the
selected regression on all outer-training sources. Inner-validation
clips may have contributed without labels to encoder fitting.

The readout summarizes five bilateral pairs using feature means (960
inputs), expanded with standard deviations, absolute adjacent-block
changes and ten support fractions (2,890 inputs). We pool held-out
predictions within each seed, give each source video equal total weight,
and average the five \(R^2\) scores. The 2,000 paired source-bootstrap
resamples retain each video's clips and seed predictions together. Their
95\% intervals condition on fitted models and splits, excluding
retraining, partition and development-selection uncertainty. Repeated
seeds and mask draws reuse the same 93 videos.

\section{Feature matching and movement estimates give different
results}\label{feature-matching-and-movement-estimates-give-different-results}

With the expanded summary, initial features yield mean \(R^2=0.223\),
compared with \(0.101\)--\(0.114\) for trained teachers (Figure 2).
Every training condition scores below its initialization in every seed;
all five paired intervals are below zero (Table 1). Mean absolute error
increases by 0.0020--0.0025 in the dimensionless contrast, from 0.0415
to 0.0436--0.0440. Final online encoders also score below their initial
weights.

\begin{center}
\begin{tabular}{lrrl}
\toprule
Training mask & Mean $R^2$ & $\Delta R^2$ & 95\% source interval \\
\midrule
Random (motion) & 0.114 & -0.109 & [-0.170, -0.041] \\
MAMP-style & 0.113 & -0.110 & [-0.165, -0.051] \\
Motion mixture & 0.114 & -0.108 & [-0.171, -0.042] \\
Random (region) & 0.109 & -0.113 & [-0.166, -0.060] \\
Connected region & 0.101 & -0.122 & [-0.183, -0.055] \\
\bottomrule
\end{tabular}
\end{center}

Table 1: trained teacher minus matched initialization, using the
expanded summary. Exploratory intervals resample source videos with
fitted models fixed; differences precede rounding.

Expanding the initial summary raises \(R^2\) from 0.071 to 0.223
(\(+0.152\)). This combines several changes: more regression inputs,
temporal variation, absolute feature increments and observation support.
Standard deviations discard order, and absolute increments discard
direction, so the gain cannot be attributed specifically to temporal
order. The initial control already includes a pretrained pose detector,
preparation, anatomical identities and supervised regression.

\begin{figure}[htbp]
\centering
\includegraphics[width=\linewidth]{assets/v8/figures/learning_results.pdf}
\caption{Different outcomes from the same fitted models. A: movement estimates from the matched initial encoder and trained teachers with the expanded summary; dots are seed scores and ticks their mean. B: correct-clip feature error is lower in 375/375 trained checks and 33/75 distinct initial checks. These checks reuse models and videos.}
\end{figure}

The predictor diagnostic compares source-weighted mean squared feature
error against correct-clip targets and valid targets from another
source, holding clips and hidden positions fixed. Correct targets have
lower error in 375/375 trained checks and 33/75 initial checks. These
totals cover 125 trained models and 25 deduplicated initial controls,
each under left-leg, right-leg and scattered gaps. Pose, viewpoint and
missingness shared by visible context and full-clip teacher features
could support this correspondence alongside motion. Changing teachers
and target distributions also prevent raw training losses from providing
a common ranking of representation quality.

Mask construction works in the inspection draws: motion weighting
selects higher-motion targets, and regions reduce targets with both
immediate temporal neighbors visible from 69.9\% to zero. Motion masks
hide approximately 17\% of valid tokens and regions 9.9\%, requiring
distinct random references. Their movement-readout intervals cross zero,
leaving mask benefits uncertain without establishing equivalence
(Appendix A).

Input preparation also changes the observable. On 623 clips from 92
sources with finite values on both paths, a retained calculation
comparison gives 70.4\% source-weighted sign agreement and
\(R^2=0.218\). This disagreement between measurement paths neither
bounds achievable model accuracy nor isolates the effect of an
individual preparation step.

\section{Implications and a future-motion
test}\label{implications-and-a-future-motion-test}

For this gait measurement, greater consistency in hidden-feature
matching accompanies poorer movement estimates under the expanded
readout. Another readout may access information that ridge regression
leaves unused: 49/125 teacher fits and 96/125 online fits select the
largest tested penalty, 10,000. Wider training-only penalty selection,
support-only controls and summaries of matched dimension would help
distinguish readout limitations from changes in the features.
Source-video holdout leaves participant independence, chronological
generalization and an untouched development cohort unresolved; recording
and pose-estimation cues also remain possible explanations.

The next experiment would predict future coordinates from an observed
prefix, with preparation, normalization and mask selection restricted to
the past. Training evidence for past-feature readout and
predicted-future-feature decoding is currently synthetic; real-data
horizon counts establish eligibility only. A real-data comparison should
use common valid endpoints and physical-time horizons, with
last-position, constant-velocity, direct-pose, initial-encoder and
mismatched-training-future controls. Decoders and their settings must
use training sources only, and future-coordinate or visibility changes
must leave the prefix and forecast unchanged. Observed future-teacher
features provide a decoder diagnostic; only features predicted from the
past test forecasting. Several horizons from one prefix still require a
separate recursive-rollout evaluation.

The present result concerns one small gait model, cohort and readout.
Foundation-model scale, broad transfer, persistent temporal state and a
working simulator require evidence beyond this study.

\section{Related work}\label{related-work}

Prediction targets and available context influence which movement
differences a representation retains. Skeleton methods predict learned
features \citep{sjepa} or explicit motion \citep{mamp}, while
connected-body masks reduce nearby cues \citep{slim}. Their
action-recognition evaluations assess suitability for labels. Video
anticipation and action-conditioned planning \citep{vjepa2} ask
different questions about behavior, and separating invariant from
transformation-sensitive features \citep{seqjepa} makes this dependence
on the intended task explicit. Our paired initial-versus-trained
comparison examines a particular observed gait quantity within that
research agenda. The mask variants borrow sampling ideas from related
methods, whose complete systems are not experimentally compared here.

\clearpage
\bibliographystyle{unsrtnat}
\bibliography{assets/v8/references}
\clearpage
\appendix
\section{Implementation and supporting
checks}\label{implementation-and-supporting-checks}

Of 666 selected annotations, 642 match pose archives and 625 pass
quality checks. Archive coordinates are \(p_x=(x_0+uW_c)/W\),
\(p_y=(y_0+vH_c)/H\) and \(p_z=dW_c/W\), for crop origin \((x_0,y_0)\),
crop size \(W_c,H_c\), frame size \(W,H\) and pose estimates \(u,v,d\).
Body scale is the median valid shoulder/hip width, with unit fallback.
The target requires observed hips; the encoder substitutes the clip's
median valid pelvis midpoint when hips are missing (zero if none
exists). Gap filling is linear between observed endpoints, without
boundary extrapolation.

Training uses width 96, encoder/predictor depths 4/2, four heads, batch
20, AdamW learning rate 0.001, weight decay 0.05, gradient clipping at 1
and teacher EMA 0.999. Geometric views rotate up to 8 degrees in the
horizontal-depth plane and translate up to 0.03 in the image plane; the
teacher receives unaugmented input. Teacher/predictor temperatures are
0.06/0.10, with center EMA 0.9. The regularizer's two-layer
96-dimensional projector uses GELU; internal
invariance/variance/covariance weights are 25/25/1. The motion target
count is half the batch's smallest valid twelve-joint gait pool, rounded
down, sampled across 33 joints. Motion weights use prepared-index
displacements: MAMP-style temperature is 0.8; the mixture combines 25\%
uniform with 75\% motion weights capped at the 95th positive-score
percentile.

Summary means and standard deviations use common paired validity;
increments require consecutive common valid blocks. Missing support
produces zero values. Each pair contributes six 96-vectors (left/right
differences and sums) and two support fractions, totaling 2,890. Ridge
candidates are \(0.01,0.1,1,10,100,1000,10000\). For \(V\) videos with
\(n_v\) clips, \(w_{vi}=1/(Vn_v)\) and \(\bar y_w=\sum w_{vi}y_{vi}\)
define

\[R_w^2=1-\frac{\sum_{vi}w_{vi}(y_{vi}-\hat y_{vi})^2}{\sum_{vi}w_{vi}(y_{vi}-\bar y_w)^2}.\]

The test mean is a metric reference; a deployable constant prediction
must be fitted on training sources. Test folds contain 19/19/19/18/18
videos and 189/182/72/77/105 clips. Retained audits check archive
integrity and split adherence. Missing person IDs and acquisition
records leave cross-upload/person deduplication and exact cohort
reconstruction unresolved.

MAMP-style, motion-mixture and region masks minus their own random
references give \(\Delta R^2=-0.0008\), \(+0.0005\) and \(-0.009\), with
intervals \([-0.020,0.015]\), \([-0.015,0.015]\) and \([-0.033,0.018]\).
A limited 264-value coordinate/validity baseline reaches \(R^2=0.035\).
Grid intervals are exploratory percentiles (bootstrap seed 812), without
familywise or development-selection adjustment. The supplement retains
aggregate seed scores and recorded intervals. Raw predictions and
checkpoints were unavailable for this revision, which checks aggregate
consistency rather than rerunning inference or source bootstraps.

Reflection augmentation has separate real-data summary evidence (mirror
probability 0.5, scheduled EMA, compact readout, four inner folds):
relative to training without mirrors, normalized teacher-token
reflection error decreases by \(\Delta q=-0.008\) \([-0.010,-0.007]\),
while \(\Delta R^2=0.004\) \([-0.006,0.013]\) leaves the movement
benefit uncertain. The comparison swaps joint identities while leaving
feature channels unchanged. An explicit reflection loss has synthetic
evidence only and adds two full encoder passes. The construction
\(g(x)=[h(x)-h(Mx)]/2\) guarantees sign reversal independently of
accuracy.

Clinical asymmetry motivates this observable \citep{stroke}, but pair
contrasts can cancel and affected-side labels and clinical validation
are absent. Videos and derived poses are excluded from the supplement;
public annotation access does not authorize their redistribution.

\textbf{AI assistance.} An AI assistant supported author-side evidence
reconciliation, literature checking, revision, figure programming and
adversarial checks. No new training was performed. Authors remain
responsible for the text, figures, claims and references.

\end{document}
